\documentclass[11pt, a4paper]{article}
\usepackage{amsthm}
\usepackage{amsmath}
\usepackage{amssymb}
\usepackage{geometry}
\geometry{
    margin=1in,
    headheight=20pt,
    footskip=10mm
}
\title{9}
\author{Nagusa, G}
\date{\today}

\begin{document}
\maketitle
%Let $F = \mathbb{R} \times \mathbb{R}$ with the identity element$0_F = (0,0)$ under $\oplus$.\\
\section*{Problem}
Define two binary operations $\oplus$ and $\odot$ on the Cartesian product $\mathbb{R}\times \mathbb{R}$ as follows:\\
for $(a,b), (c,d) \in \mathbb{R}\times \mathbb{R}$
\begin{align*}
    (a,b) \oplus (c,d) &:= (a+c, b+d) \\
    (a,b) \odot (c,d) &:= (ac-bd, ad+bc)
\end{align*}
It is well known that $\langle \mathbb{R} \times \mathbb{R}, \oplus \rangle$ is an abelian group.
Prove that $\langle \mathbb{R}\times \mathbb{R}, \oplus, \odot \rangle$ is a field.

\section*{Proof}
Let $F = \mathbb{R} \times \mathbb{R}$ with the identity element $0_F = (0,0)$ under $\oplus$.\\
Since it is given that $\langle F, \oplus \rangle$ is an abelian group,
we only need to verify the following to prove that $F$ is a field:
\begin{enumerate}
    \item $\odot$ is commutative
    \item $\odot$ is associative
    \item $\odot$ and $\oplus$ satisfy the distributive law
    \item Existence of the unity $1_F \neq 0_F$
    \item Every nonzero element is a unit
\end{enumerate}
Let $x = (a,b), y = (c,d)$, and $z = (e,f) \in F$.\\
Note that multiplication and addition in $\mathbb{R}$ are commutative and associative.

\subsection*{1. $\odot$ is commutative}
\begin{align*}
    x \odot y &= (a,b) \odot (c,d) \\
    &= (ac-bd, ad+bc) \\
    &= (ca-db, cb+da) \\
    &= (c,d) \odot (a,b) = y \odot x
\end{align*}

\subsection*{2. $\odot$ is associative}
\begin{align*}
    (x \odot y) \odot z &= (ac-bd, ad+bc) \odot (e,f) \\
    &= ((ac-bd)e - (ad+bc)f, (ac-bd)f + (ad+bc)e) \\
    &= (ace - bde - adf - bcf, acf - bdf + ade + bce)\\
    &= (ace - adf - bcf - bde, acf + ade + bce - bdf)\\
    &= (a(ce-df) - b(cf+de), a(cf+de) + b(ce-df)) \\
    &= (a,b) \odot (ce-df, cf+de) \\
    &= x \odot (y \odot z)
\end{align*}

\subsection*{3. $\odot$ and $\oplus$ satisfy the distributive law}
\begin{align*}
    x \odot (y \oplus z) &= (a,b) \odot (c+e, d+f) \\
    &= (a(c+e) - b(d+f), a(d+f) + b(c+e)) \\
    &= (ac+ae - bd-bf, ad+af + bc+be) \\
    &= (ac-bd + ae-bf, ad+bc + af+be) \\
    &= (ac-bd, ad+bc) \oplus (ae-bf, af+be) \\
    &= (x \odot y) \oplus (x \odot z)
\end{align*}
(Since $\odot$ is commutative, the right distributive law also follows.)

\subsection*{4. Existence of the unity $1_F \neq 0_F$}
Consider the element $(1,0) \in F$. \\
For any $(a,b) \in F$, 
\[
    (a,b) \odot (1,0) = (a\cdot 1 - b \cdot 0, a \cdot 0 + b \cdot 1) = (a,b)
\]
Since $0_F = (0,0)$, we have $(1,0) \neq 0_F$. \\
Thus, $(1,0) = 1_F$ is the unity of $F$.

\subsection*{5. Every nonzero element is a unit}
Let $x = (a,b) \in F$ be a nonzero element, which means $(a,b) \neq (0,0)$.\\
Then $a^2 + b^2 \neq 0$ since $a,b \in \mathbb{R}$.\\
Consider the element $x^{-1} = \left( \frac{a}{a^2+b^2}, \frac{-b}{a^2+b^2} \right) \in F$.
\begin{align*}
    x\odot x^{-1} &= (a,b) \odot \left( \frac{a}{a^2+b^2}, \frac{-b}{a^2+b^2} \right)\\
    &= \left( a\left(\frac{a}{a^2+b^2}\right) - b\left(\frac{-b}{a^2+b^2}\right), a\left(\frac{-b}{a^2+b^2}\right) + b\left(\frac{a}{a^2+b^2}\right) \right) \\
    &= \left( \frac{a^2+b^2}{a^2+b^2}, \frac{-ab+ab}{a^2+b^2} \right) \\
    &= (1,0)\\
    &= 1_F
\end{align*}
(Since $\odot$ is commutative, $x^{-1} \odot x = 1_F$ also holds.)\\
We have that $x^{-1}$ is a multiplicative inverse of $x$.\\
Thus, every nonzero element is a unit.

\subsection*{$\therefore$ $\langle \mathbb{R}\times \mathbb{R}, \oplus, \odot \rangle$ is a field.}
\end{document}